% First draft, 2026-10-29. Build with LuaLaTeX (Japanese via luatexja, IPA via a font that has it): see README.md.
\documentclass[11pt,a4paper]{article}
\usepackage[margin=25mm]{geometry}
\usepackage{luatexja-fontspec}
\setmainfont{Charis SIL}          % IPA, including the pitch-accent letters U+A71B / U+A71C
\setmonofont{DejaVu Sans Mono}
\setmainjfont{Noto Serif CJK JP}
\usepackage{booktabs}
\usepackage{tabularx}
\usepackage{amsmath}
\usepackage{xcolor}
\usepackage[hidelinks]{hyperref}
\usepackage[numbers,sort&compress]{natbib}

% YomiScript and IPA examples: copied exactly from the spec and the result files
\newcommand{\ys}[1]{\mbox{\texttt{#1}}}
\newcommand{\provisional}{\textsuperscript{\textdagger}}
\newcommand{\todo}[1]{\textcolor{red}{[TODO: #1]}}

\title{The Author Knows the Reading:\\
LLM-Written Pronunciation for Small Browser Text-to-Speech, with Japanese Homographs}
\author{\todo{authors and affiliations}}
\date{Draft, 29 October 2026}

\begin{document}
\maketitle

\begin{abstract}
Text-to-speech (TTS) front ends choose a reading for every word before any audio is made, and in Japanese that choice
is often ambiguous: the same written form has several readings (homographs such as 一日, 上手, 人気), numbers and dates
change their sound with the counter, and Latin-script words must be rendered in Japanese phonology. Dictionary-based
front ends such as OpenJTalk pick the most frequent reading and do not look at context. We observe that when speech is
produced for a large language model (LLM), the LLM that wrote the sentence already knows which reading it meant. We
let that LLM write the pronunciation itself, in YomiScript, an IPA-based, language-neutral notation that separates the
\emph{Yomite} (who decides how text is read) from the \emph{Katarite} (the synthesizer that turns the resolved
pronunciation into sound). The TTS then needs no front end of its own: only a small acoustic model, which runs in a
web browser. On a 70-sentence hard set of Japanese homographs, acronyms, numbers and dates, LLM-written YomiScript was
read correctly in 100\% (Codex) and 98.6\% (Claude) of sentences, against 78.6\% for pyopenjtalk and 32.9\% for the
original front end of Kokoro-82M; on the 35 homograph sentences alone, 100\% and 97.1\% against 60.0\% and 51.4\%.
We describe a Katarite for Kokoro that reproduces the original Kokoro exactly when given its own front end's output,
a second Katarite for a text-reading engine (Supertonic 3) that receives kana only for ambiguous words, and an
optional original-text layer that lets the same YomiScript drive both and show the text being spoken. These are
preliminary results with provisional gold readings, small sets and one run per LLM; we list what is needed to confirm them.
\end{abstract}

\section{Introduction}

Speech interfaces to LLMs usually send the generated text to a cloud speech API. Moving synthesis onto the device,
into the browser tab that shows the chat, removes the network round trip and the per-character cost, and small
open models such as Kokoro-82M~\citep{kokoro82m} now make this practical. What does not move easily is the front end:
the grapheme-to-phoneme (G2P) step that decides how each word is read. For English it is a pronunciation dictionary
plus rules; for Japanese it is a morphological analyser with a large dictionary. Kokoro's own Japanese front end
(misaki~\citep{misaki}, built on MeCab~\citep{kudo2004mecab} with UniDic 3.1.0~\citep{den2008unidic}) needs a
connection-cost table of 481\,MB (413\,MB compressed with gzip) and a 243\,MB dictionary; it cannot be shipped to a
browser.

Even where a front end fits, it reads without understanding. In Japanese the problem is common: 一日 is
\emph{tsuitachi} (the first day of the month) or \emph{ichinichi} (one day), 上手 is \emph{jōzu} (skilful),
\emph{kamite} (stage left) or \emph{uwate} (superior), 人気 is \emph{ninki} (popularity) or \emph{hitoke} (signs of
people). OpenJTalk~\citep{openjtalk} and its Python wrapper pyopenjtalk~\citep{pyopenjtalk} choose one reading per
written form from the dictionary; on our hand-made hard set pyopenjtalk reads 60.0\% of the homograph sentences correctly
under the scoring of Section~\ref{sec:setup} (51.4\% under the strict, spelling-sensitive scoring of an earlier baseline run).

The observation behind this work is simple: in an LLM application, the text to be spoken was written by an LLM, and
that LLM knows what it meant. If it writes the pronunciation as well as the text, the hardest part of the front end
is done by the component that has the context. What remains for the TTS is to turn a resolved pronunciation into
sound, which a small acoustic model can do in the browser.

Our contributions:
\begin{itemize}
  \item \textbf{YomiScript}, a pronunciation notation in IPA~\citep{ipa1999handbook} with accent, stress and tone
  marks, phrase boundaries, optional pitch values and an optional original-text layer, which separates the
  \emph{Yomite} (text to YomiScript: an LLM, rules, a dedicated model or a person) from the \emph{Katarite}
  (YomiScript to sound) (Section~\ref{sec:yomiscript}).
  \item A \textbf{Kokoro Katarite} that, given YomiScript written from Kokoro's own front end, hands the model exactly
  the phoneme string and ids of the original pipeline, with identical durations and F0, and audio within the
  original's own run-to-run variation (Section~\ref{sec:parity}).
  \item An evaluation of \textbf{LLMs as the Yomite} for Kokoro on Japanese and English, against the original front
  end and pyopenjtalk, on readings, faithfulness of the text layer, intelligibility and cost (Section~\ref{sec:results}).
  \item A \textbf{second Katarite} for an engine that reads characters rather than IPA (Supertonic~3~\citep{supertonic3}),
  which uses the text layer and gives kana only for words whose written form is ambiguous.
\end{itemize}

\section{Related work}

\paragraph{Japanese front ends.} OpenJTalk~\citep{openjtalk} (and pyopenjtalk~\citep{pyopenjtalk}) analyses text with
MeCab and a NAIST-JDIC-based dictionary and produces readings, accent phrases and accent nuclei; it is the de facto
front end of open Japanese TTS. misaki~\citep{misaki}, the G2P of Kokoro, reads Japanese with UniDic through a
romanisation-based path. MARINE~\citep{park2022marine} estimates accent with multi-task learning. zenz~\citep{zenz}
models the reverse direction (kana-kanji conversion) with a small GPT-2 model. All choose readings from a dictionary
or a model trained on dictionary output; none is given the writer's intent.

\paragraph{Neural G2P, homographs and text normalisation.} Multilingual neural G2P such as ByT5-based
models~\citep{zhu2022byt5g2p} learns pronunciations from lexicons. Homograph disambiguation with contextual
embeddings has been studied for English TTS~\citep{nicolis2021homograph}, and text normalisation (numbers, dates,
symbols into words) is a long-standing hard problem~\citep{sproat2016normalization}. LLMs have been tried for G2P and
normalisation \todo{cite LLM-based G2P and normalisation work, e.g. benchmarks of LLM G2P; verify}. We differ in who
does it: not a separate model reading someone else's text, but the author of the text.

\paragraph{Small and browser TTS.} Kokoro-82M~\citep{kokoro82m} is a StyleTTS~2-style~\citep{li2023styletts2} model with
an iSTFTNet~\citep{kaneko2022istftnet} decoder that takes phonemes. Piper~\citep{piper}, based on VITS~\citep{kim2021vits},
runs in the browser through ONNX Runtime~\citep{onnxruntime}; piper-plus~\citep{piperplus} adds Japanese with OpenJTalk
compiled to WebAssembly. kokoro-js~\citep{kokorojs} runs Kokoro in the browser with English voices only.
SupertonicTTS~\citep{kim2025supertonic} reads raw characters (no G2P), learns text-speech alignment with flow
matching~\citep{lipman2023flow} and predicts only the utterance duration; Supertonic~3~\citep{supertonic3} covers 31
languages including Japanese. Phonon~\citep{phonon} is a 31M/40M on-device model with a WebAssembly build, without
Japanese at launch. The Web Speech API~\citep{webspeech} exposes the operating system's voices. In all of them the
reading is decided inside the TTS, by a dictionary front end or by the acoustic model itself.

\section{YomiScript}\label{sec:yomiscript}

\paragraph{Roles.} A \emph{Yomite} (読み手, the reader) turns text into YomiScript; a \emph{Katarite} (語り手, the
narrator) turns YomiScript into sound. Lexicons and reading hints are inputs to the Yomite and are not part of
YomiScript, which holds only resolved pronunciation.

\paragraph{Notation.} An utterance is a BCP~47 language tag and the pronunciation between slashes, ending in a fall
(\ys{↘}) or rise (\ys{↗}). Every language is written in IPA, Japanese included, mora by mora, with \ys{ꜛ} after the
mora where the pitch rises and \ys{ꜜ} after the accent nucleus; English uses \ys{ˈ}/\ys{ˌ}. \ys{|} is an accent-phrase
boundary and \ys{‖} a breath group with a pause. Optional pitch values \ys{<f0:v>} give a target in semitones relative
to the speaker's register for the unit before them; a Katarite that has its own prosody model may ignore marks and fill
pitch itself. Examples from the specification:

\begin{quote}
\ys{ja /kʲoꜜːwa | joꜜi | teꜛŋkidesɯ ↘/}\\
\ys{en-US /ˈaɪ ˈɹɛd ðə ˈbʊk | ˈjɛstɚˌdeɪ ↘/}
\end{quote}

\paragraph{Original text.} \ys{<t:…>} closes a span: the units since the previous \ys{<t:…>} are the reading of that
text, and a \ys{<t:…>} with no units before it is text that is not spoken (punctuation). Joining the values gives back
the text. The layer never changes the pronunciation: the parser removes it before reading the IPA, and Kokoro's input
is identical with or without it (checked on the 48 fixtures of Section~\ref{sec:parity}). It serves two purposes:
showing the text while it is spoken (a chat client highlighting the current word), and Katarites whose engine reads
characters. An LLM output from our hard set, for 一日から三日まで休みです (``closed from the 1st to the 3rd''):

\begin{quote}
\ys{ja /tsɯitatɕi<t:一日> kaɾa<t:から> mikka<t:三日> made<t:まで> jasɯmi<t:休み> desɯ<t:です。> ↘/}
\end{quote}

and for 舞台の上手から役者が登場した (``an actor entered from stage left''), where 上手 is \emph{kamite}:

\begin{quote}
\ys{ja /bɯtai<t:舞台> no<t:の> kamite<t:上手> kaɾa<t:から> jakɯɕa<t:役者> ɡa<t:が> toːdʑoː<t:登場> ɕita<t:した。> ↘/}
\end{quote}

\paragraph{Katarites.} The \emph{Kokoro Katarite} maps YomiScript IPA to Kokoro's phoneme inventory and runs the
model's three ONNX parts (text encoder with duration prediction, F0/energy prediction, decoder) in the browser or in
Node. The \emph{Supertonic Katarite} gives the engine text: each span as written, except spans whose written form (or
a 2--8 character form inside it) has more than one reading in a list of 12{,}875 forms derived from
IPAdic~\citep{ipadic}, or that contain numbers, Latin letters, symbols or kanji numerals; those are given as kana from
the YomiScript reading (は as わ, long vowels as ー). The list and the rule belong to the Katarite; the Yomite writes
only YomiScript.

\section{Experimental setup}\label{sec:setup}

\paragraph{Sets.} (1)~\emph{ja}: 24 Japanese texts written to probe Kokoro's front end (punctuation, questions, numbers,
loanwords, Latin script, a long passage); (2)~\emph{en}: 22 English texts of the same kind; (3)~\emph{hard}: 70
Japanese sentences, 35 homograph sentences, 12 with acronyms, 12 with numbers and 11 with dates, with gold readings in
katakana. The gold readings are provisional (Section~\ref{sec:limits}).

\paragraph{Yomites.} Two LLMs through their command-line agents with all tools disabled: Claude Code
(\texttt{claude-opus-5-5}) and Codex (\texttt{gpt-6.1-sol}), as reported by the CLIs. One call per set, with a prompt
giving the subset of YomiScript Kokoro needs (IPA, word spaces, \ys{‖}, \ys{↘}/\ys{↗}, no pitch values or accent marks,
since Kokoro predicts prosody), the Japanese mora-to-IPA table of the specification, the text layer, and five examples
written by Kokoro's own front end on sentences outside the test sets. Baselines: Kokoro's original front end (misaki,
via KPipeline) and pyopenjtalk~0.4.1 (the same OpenJTalk and NAIST-JDIC front end that piper-plus compiles to
WebAssembly; we use it as piper-plus's stand-in for readings).

\paragraph{Prompt rounds.} The prompt was revised twice, for notation only: the first round gave too little of it (the
LLMs wrote phonemic \ys{ni} for に where the specification writes \ys{ɲi}, and several utterances per item); the second
added the specification's mora table and the one-utterance rule; the third added the text layer. The instruction about
readings (``each word read the way it is meant in its context'') did not change. All numbers are from the third round
unless stated otherwise.

\paragraph{Scoring.} Every output goes through the Kokoro Katarite's planner to the phoneme string Kokoro is given. On
\emph{ja} and \emph{en} we compare with the original front end; on \emph{hard} with the gold. A reading counts as
\emph{the same} when the phonemes match after removing word spaces and punctuation and merging spellings that do not
change the reading: long-vowel spelling (イイ/イー), Kokoro's ɨ for a devoiced ɯ, the allophone written for a moraic
nasal, and h/n before i written without their palatal allophone (for English: stress marks, Kokoro's reduced vowels and
flap). Intelligibility is measured by transcribing the audio with Whisper large-v3-turbo~\citep{radford2023whisper},
converting the transcript to katakana with pyopenjtalk, and computing the kana error rate against what the system was
meant to say (English: word error rate against the text).

\section{Results}\label{sec:results}

\subsection{The Kokoro Katarite reproduces the original}\label{sec:parity}

Given YomiScript written from Kokoro's own front end, the Katarite hands Kokoro the same input as the original
pipeline (Table~\ref{tab:parity}). With fp32 weights the audio differs from the original no more than the original
differs from itself (its decoder draws random numbers); the fp16-weight model used in the browser is close but not
identical. This makes the comparisons below about readings, not about a different synthesizer.

\begin{table}[t]
\centering\small
\caption{Kokoro Katarite vs the original Kokoro pipeline (kokoro 0.9.4 KPipeline), given the original front end's
pronunciation as YomiScript. Correlation of the waveform with the original's; the baseline is a second run of the original.}
\label{tab:parity}
\begin{tabular}{lcccc}
\toprule
& \multicolumn{2}{c}{English (23 chunks)} & \multicolumn{2}{c}{Japanese (50 chunks)}\\
\cmidrule(lr){2-3}\cmidrule(lr){4-5}
Phoneme string, ids, voice index, durations & \multicolumn{2}{c}{23/23 identical} & \multicolumn{2}{c}{50/50 identical}\\
F0 (RMSE) & \multicolumn{2}{c}{0.000 st} & \multicolumn{2}{c}{0.000 st}\\
\midrule
Audio correlation & mean & min & mean & min\\
Original, second run & 0.9943 & 0.9924 & 0.9918 & 0.9762\\
Katarite, fp32 & 0.9944 & 0.9888 & 0.9914 & 0.9773\\
Katarite, fp16 weights (browser) & 0.9854 & 0.8545 & 0.9833 & 0.7980\\
\bottomrule
\end{tabular}
\end{table}

\subsection{Readings}

\begin{table}[t]
\centering\small
\caption{Sentences read correctly on the hard set (gold\provisional). One run per LLM.}
\label{tab:hard}
\begin{tabular}{lccccc}
\toprule
System & all (70) & homograph (35) & acronym (12) & number (12) & date (11)\\
\midrule
LLM Yomite: Codex & \textbf{100.0\%} & \textbf{100.0\%} & 100.0\% & 100.0\% & 100.0\%\\
LLM Yomite: Claude & 98.6\% & 97.1\% & 100.0\% & 100.0\% & 100.0\%\\
pyopenjtalk (piper-plus front end) & 78.6\% & 60.0\% & 100.0\% & 91.7\% & 100.0\%\\
Kokoro's original front end (misaki) & 32.9\% & 51.4\% & 0.0\% & 41.7\% & 0.0\%\\
\bottomrule
\end{tabular}
\end{table}

Table~\ref{tab:hard} shows the main result. Both LLMs read the homographs from context; pyopenjtalk reads 21 of 35 and
Kokoro's front end 18 of 35. Kokoro's front end also leaves Latin-script acronyms (NHK, SAP) as they are and misreads
dates (24日 as \emph{nijūyon-nichi} rather than \emph{nijūyokka}). Claude's one miss (今日の日本経済を論じる) is a
notation slip rather than a misreading: it wrote \ys{keːzai} where the specification writes ザ as \ys{dza}; the reading
of the homograph 今日 (\emph{konnichi}) is correct. In the first prompt round, when the LLMs had not been given the mora
table, the same sets scored 55.7\% and 51.4\% on the hard set, almost entirely because of such notation mismatches; we
report this to make clear that the notation, not only the model, determines the score.

\begin{table}[t]
\centering\small
\caption{Agreement with Kokoro's original front end on the parity sets (same reading after merging spellings; Section~\ref{sec:setup}).}
\label{tab:paritysets}
\begin{tabular}{lcccc}
\toprule
& \multicolumn{2}{c}{Japanese (24 texts)} & \multicolumn{2}{c}{English (22 texts)}\\
\cmidrule(lr){2-3}\cmidrule(lr){4-5}
System & same reading & phoneme error & same reading & phoneme error\\
\midrule
LLM Yomite: Claude & 62.5\% & 9.22\% & 31.8\% & 10.36\%\\
LLM Yomite: Codex & 62.5\% & 8.97\% & 40.9\% & 9.39\%\\
pyopenjtalk & 62.5\% & 10.92\% & -- & --\\
\bottomrule
\end{tabular}
\end{table}

Agreement with Kokoro's own front end (Table~\ref{tab:paritysets}) is not a quality measure: in Japanese most
differences are the original's misreadings (清水寺 as \emph{shimizu-tera} for \emph{kiyomizu-dera}, 9時 as
\emph{kyū-ji} for \emph{ku-ji}, iPhone left in Latin letters). In English the differences are mostly phonetic detail
(stress placement, Kokoro's reduced-vowel symbols), which the model does hear; we leave English to listening tests.

\subsection{Text layer and cost}

Joining the \ys{<t:…>} values gave back the input exactly (whitespace ignored) in all 232 outputs (2 LLMs $\times$ 116
texts). The layer roughly doubles the output: the YomiScript is 5.6--6.1 times the characters of the Japanese text
(2.7--3.2 without it) and 3.1 times for English; Claude wrote 10{,}692, 8{,}424 and 11{,}111 output tokens for the
\emph{ja}, \emph{en} and \emph{hard} sets. The calls took 31--112\,s per set (\emph{not} a latency measure: one call
for a whole set).

\subsection{Intelligibility and the two Katarites}

\begin{table}[t]
\centering\small
\caption{Kana error rate of Whisper transcripts on the hard set (lower is better), against the gold reading. A proxy:
Whisper and pyopenjtalk add their own errors, equally for all systems.}
\label{tab:asr}
\begin{tabular}{lcc}
\toprule
System & all (70) & homographs (35)\\
\midrule
Claude YomiScript $\to$ Kokoro & 3.7\% & 6.6\%\\
Codex YomiScript $\to$ Kokoro & 3.9\% & 7.0\%\\
Codex YomiScript $\to$ Supertonic Katarite (kana for ambiguous words) & 5.6\% & 8.2\%\\
Claude YomiScript $\to$ Supertonic Katarite (kana for ambiguous words) & 7.6\% & 9.2\%\\
macOS system voice (Kyoko) & 5.5\% & 10.1\%\\
Kokoro with its original front end & 10.3\% & 9.6\%\\
Supertonic 3 on the plain text & 13.8\% & 14.4\%\\
\bottomrule
\end{tabular}
\end{table}

With the same YomiScript, Kokoro was the more intelligible Katarite (Table~\ref{tab:asr}). Supertonic~3 on the plain
text had the highest error, and giving it kana for ambiguous words, chosen by its Katarite from the YomiScript and the
text layer, cut the error on the hard set by about half. The rule left 3 of 35 homographs as written (市役所, 朝市,
何本), all in words that fix the reading. In informal listening by the authors, Supertonic sounded more natural than
Kokoro when its reading was right; Whisper cannot measure this, and we have no listening test yet. On English the word
error rates were close (2.0--3.2\%).

\subsection{Context: other Yomites we tried}

In an earlier experiment (7 Japanese items) the LLMs wrote complete YomiScript with accent marks and pitch values from
plain text; their accent-nucleus F1 against a rule-based reference was 0.82 (Claude) and 0.60 (Codex), and adding pitch
values to YomiScript whose marks were given was easy for both. Accent from plain text is the weaker part, so here the
Katarite's engine predicts prosody. A dedicated 7.7M-parameter reading model (Yomi-0) trained on
300k sentences labelled by pyopenjtalk stayed below its teacher (JSUT~\citep{sonobe2017jsut,jsutlabel} sentence
accuracy 68.2\% vs 72.5\%; hard set 57.1\% vs 71.4\%): distillation inherits the teacher's context-free readings, which
is the problem this work avoids.

\section{Discussion}

\paragraph{Why it works.} The author of the text has the context, and in an LLM application the author is a model that
can also write IPA. A dictionary front end must guess; the author does not.

\paragraph{What it buys for the browser.} No front end ships: the browser needs only the Katarite (Kokoro with fp16
weights, 163\,MB plus a 0.5\,MB voice; Supertonic~3, 385\,MB). Compared with shipping OpenJTalk (piper-plus: about
19\,MB of WebAssembly, by its documentation), the gain is not size but readings.

\paragraph{Delivery.} The same idea fits a Model Context Protocol~\citep{mcp} server or a chat-client plugin: the
calling LLM writes YomiScript as the argument of a \texttt{speak} tool, and the client can show the text from the text
layer while it is spoken. The TTS has no LLM and no front end of its own.

\paragraph{Costs.} The calling LLM writes several times more characters than the text, which delays the first audio and
costs tokens; streaming per sentence and a lighter notation (for example omitting the text layer when not displayed)
are ways to reduce it.

\section{Limitations}\label{sec:limits}

The following weaken the claims above, and Section~\ref{sec:results} should be read with them.
\begin{itemize}
  \item \textbf{Provisional gold (weakens Table~\ref{tab:hard}).} The hard set and its gold readings were written by
  the project, marked provisional and not reviewed by independent annotators; LLM assistance in writing them cannot be
  ruled out, which could favour the LLMs. Several items admit two readings.
  \item \textbf{Small sets (weakens all tables).} 35 homograph sentences and 70 hard sentences; one wrong answer moves
  the homograph score by 2.9 points. No confidence intervals are reported.
  \item \textbf{One run per LLM (weakens Tables~\ref{tab:hard}--\ref{tab:asr}).} LLM outputs vary between calls; one
  Codex call returned malformed JSON and was repeated.
  \item \textbf{Prompt revisions.} The prompt was revised twice after seeing results, for notation; the reading
  instruction did not change, but the revision was informed by the test sets.
  \item \textbf{Two frontier LLMs only.} Smaller or local models, which an MCP server will also meet, were not tested.
  \item \textbf{Intelligibility is a proxy (weakens Table~\ref{tab:asr}).} ASR transcripts, read back with pyopenjtalk,
  carry both systems' errors; naturalness was not measured.
  \item \textbf{Baselines.} pyopenjtalk stands in for piper-plus; context-aware Japanese G2P models were not compared.
  \item \textbf{Long utterances.} The Kokoro Katarite does not yet split an utterance longer than Kokoro's 510 phonemes;
  the one long passage could not be spoken from a single LLM utterance.
\end{itemize}

\section{Conclusion}

When an LLM speaks, it can also say how its words are read. Writing that reading in a language-neutral notation lets a
small synthesizer in the browser speak Japanese homographs, numbers, dates and acronyms that dictionary front ends
misread, with the reading decided by the one component that knows what was meant. Our preliminary results are strong
but small; the next step is an independent, larger evaluation with human listening tests.

\bibliographystyle{plainnat}
\bibliography{references}

\end{document}
